% Presentation: Berlin Hausverwaltung Leads (Heidi Systems case study).
% Beamer source; one frame = one slide. \note{} holds what to say on each slide.
% Story: why KPIs come first -> choosing a method -> deriving the KPIs -> testing them ->
% the evidence behind them -> turning them into a score -> how review changed them ->
% the agent that fills them -> does it work -> what we learned.
\documentclass[aspectratio=169,11pt]{beamer}

\usetheme{default}
\setbeamertemplate{navigation symbols}{}
\setbeamertemplate{footline}[frame number]
\usepackage[T1]{fontenc}
\usepackage[utf8]{inputenc}
\usepackage{booktabs}

\definecolor{accent}{HTML}{0B6A70}
\definecolor{assume}{HTML}{9A4F0A}
\definecolor{bad}{HTML}{A33A3A}
\setbeamercolor{frametitle}{fg=accent}
\setbeamercolor{title}{fg=accent}
\setbeamercolor{structure}{fg=accent}
\setbeamercolor{section in toc}{fg=accent}
\newcommand{\ev}{\textcolor{accent}{\textbf{Evidence}}}
\newcommand{\as}{\textcolor{assume}{\textbf{Assumption}}}
\newcommand{\yes}{\textcolor{accent}{\textbf{pass}}}
\newcommand{\no}{\textcolor{bad}{\textbf{fail}}}
\newcommand{\partlabel}[1]{{\small\color{assume}\textsc{#1}}\par\vspace{0.3em}}

\title{Berlin Hausverwaltung Leads}
\subtitle{How I decided what makes a good lead, and built an agent to measure it}
\author{Hammad Naseer}
\date{Case study \textbullet{} AI Agent Builder \textbullet{} Heidi Systems \textbullet{} November 2026}

\begin{document}

%====================================================================
\begin{frame}
  \titlepage
  \note{This presentation follows the order I worked in. Most of it is about one decision:
  what makes a Hausverwaltung a good lead for Heidi. Everything else, including the agent,
  follows from that.}
\end{frame}

%====================================================================
\begin{frame}{The task}
  \begin{itemize}
    \item Heidi is a \textbf{metering service provider}: heating and water meters, and the
          annual heating-cost bill under the HeizkostenV
    \item Heidi's customers are \textbf{Hausverwaltungen} (property managers), not tenants
    \item The task:
      \begin{enumerate}
        \item List as many Berlin Hausverwaltungen as possible
        \item Choose qualifying KPIs and fill them in for each company
        \item Bonus: a named contact with email or LinkedIn
        \item Explain sources, tools, automation and validation
      \end{enumerate}
    \item The brief says the \textbf{choice of KPIs and the reasoning behind it} are part of
          the evaluation
  \end{itemize}
  \note{The list is the easy part. The real question in the brief is: which companies should
  Heidi call first, and why? That's a KPI question.}
\end{frame}

%====================================================================
\section{Why the KPIs come first}

\begin{frame}{Why the KPIs come first}
  \partlabel{Part 1\textbullet{} The first decision}
  The KPIs decide everything that comes after them:
  \begin{itemize}
    \item \textbf{What to collect:} which data sources are worth using
    \item \textbf{What the agent looks for} on each website
    \item \textbf{How companies are ranked}, which is what sales actually uses
  \end{itemize}
  \vspace{0.8em}
  If the KPIs are wrong, a perfect pipeline still produces a \textbf{confident, wrong ranking}.
  \vspace{0.8em}

  So I set the method for choosing KPIs \emph{before} looking at any data.
  \note{It's tempting to start scraping and then score whatever you find. But then the KPIs
  are chosen by what's easy to get, not by what predicts a good customer.}
\end{frame}

\begin{frame}{Two traps to avoid}
  \begin{columns}[T]
    \column{0.48\textwidth}
    \textbf{Trap 1: measuring what's easy}
    \begin{itemize}
      \item Google gives a rating, a district and a phone number for every company
      \item None of them says how many meters a company could bring
      \item Easy data looks like a KPI, but it doesn't predict value
    \end{itemize}
    \column{0.48\textwidth}
    \textbf{Trap 2: gut feeling}
    \begin{itemize}
      \item ``Bigger companies are better'' sounds right
      \item But how much better? Is a municipal giant a good lead?
      \item Without a method, nobody can check or challenge the choice
    \end{itemize}
  \end{columns}
  \vspace{1em}
  What I needed: a method where \textbf{every KPI traces back to a business question and
  a source}.
  \note{Example: early on I asked whether including the big municipal housing companies was
  a gut call or a rule. Without a method I couldn't answer that. With one, it became a
  measurable KPI: buyer type.}
\end{frame}

%====================================================================
\section{Choosing a method}

\begin{frame}{Which method? The options I considered}
  \partlabel{Part 2 \textbullet{} Choosing the method}
  \small
  \begin{tabular}{@{}p{3.4cm} p{5.6cm} p{3.6cm}@{}}
    \toprule
    \textbf{Method} & \textbf{What it is for} & \textbf{Decision} \\
    \midrule
    BANT / MEDDIC & Qualifying a deal already in talks: budget, authority, need, timing &
      \textcolor{bad}{Rejected:} a cold list has no budget or timing to observe \\
    Gut-feeling scorecard & Quick ranking by experience & \textcolor{bad}{Rejected:} not
      traceable or testable \\
    \textbf{Goal-Question-Metric} & Deriving metrics top-down from a goal
      (Basili et al., 1994) & \textcolor{accent}{\textbf{Chosen}} to \emph{derive} the KPIs \\
    \textbf{SMART} & Checking that a metric is usable (Doran, 1981) &
      \textcolor{accent}{\textbf{Chosen}} to \emph{filter} the KPIs \\
    \textbf{B2B segmentation} & Firmographic, technographic, behavioral &
      \textcolor{accent}{\textbf{Chosen}} to \emph{check coverage} \\
    \textbf{Weighted lead scoring} & Combining several KPIs into one rank &
      \textcolor{accent}{\textbf{Chosen}} to \emph{score} \\
    \bottomrule
  \end{tabular}
  \note{BANT and MEDDIC are the best-known sales frameworks, but they assume you're already
  talking to a buyer. For a cold list scraped from the web, budget and timing simply can't be
  observed, so they'd produce empty columns.}
\end{frame}

\begin{frame}{Why Goal-Question-Metric}
  \begin{itemize}
    \item Developed at the University of Maryland and NASA's Software Engineering Laboratory
          to stop teams from measuring \textbf{whatever is easy to count}
    \item It works \textbf{top-down}:
      \begin{enumerate}
        \item \textbf{Goal:} what do we want to achieve?
        \item \textbf{Questions:} what do we need to know to get there?
        \item \textbf{Metrics:} what measurable value answers each question?
      \end{enumerate}
    \item Every metric is justified by a question, and every question by the goal
    \item A metric that answers no question gets dropped, however easy it is to collect
  \end{itemize}
  \vspace{0.5em}
  This is exactly the protection against both traps.
  \note{The key property: you can point at any column in the sheet and ask ``why is this
  here?'', and the answer is always a question about Heidi's business.}
\end{frame}

\begin{frame}{The full chain: each step has one job}
  \begin{enumerate}
    \item \textbf{Goal-Question-Metric} \textrightarrow{} \emph{proposes} candidate KPIs
    \item \textbf{SMART} \textrightarrow{} \emph{removes} KPIs that can't be measured or found
    \item \textbf{Three tests} \textrightarrow{} \emph{ranks} the survivors:
      \begin{itemize}
        \item \textbf{Correlation:} does it actually move deal size or win chance?
        \item \textbf{Discrimination:} does it differ between companies?
        \item \textbf{Feasibility:} can we find and verify it for most companies?
      \end{itemize}
    \item \textbf{Tiers and weights} \textrightarrow{} Critical 3, Important 2, Useful 1
    \item \textbf{Label every reason} \textrightarrow{} \ev{} with a source, or \as{} stated openly
  \end{enumerate}
  \note{The tiering tests are what turn a list of reasonable ideas into a ranking. A KPI can
  be relevant and still useless if every company has the same value.}
\end{frame}

%====================================================================
\section{Deriving the KPIs}

\begin{frame}{Step 1: the goal and the five questions}
  \partlabel{Part 3 \textbullet{} Goal-Question-Metric applied}
  \textbf{Goal:} find the Berlin Hausverwaltungen that are the \textbf{most valuable} and
  \textbf{most winnable} customers for a metering service.
  \vspace{0.6em}

  \begin{tabular}{@{}p{6.6cm} p{6.4cm}@{}}
    \toprule
    \textbf{Question} & \textbf{Why it matters to Heidi} \\
    \midrule
    Q1 How much metering volume could they bring? & Heidi charges per unit per year \\
    Q2 Is the portfolio covered by HeizkostenV billing? & That's what creates the need \\
    Q3 Who signs, and how easily can it be won? & Decides the sales path \\
    Q4 How easily could they switch to Heidi? & Integration lowers the barrier \\
    Q5 Is there a reason to act now? & Timing of the sales push \\
    \bottomrule
  \end{tabular}
  \note{``Valuable'' gives Q1 and Q2, ``winnable'' gives Q3 and Q4, and Q5 is about timing.
  Each question comes from how Heidi actually makes money.}
\end{frame}

\begin{frame}{Step 2: candidate KPIs for each question}
  \small
  \begin{tabular}{@{}p{3.0cm} p{9.8cm}@{}}
    \toprule
    \textbf{Question} & \textbf{Candidates} \\
    \midrule
    Q1 Volume & Units managed \textbullet{} number of buildings \textbullet{} employees
      \textbullet{} growth (hiring, expansion) \\
    Q2 HeizkostenV & Residential share of the portfolio \\
    Q3 Who signs & Buyer type (rental manager, WEG manager, cooperative, own portfolio,
      municipal) \textbullet{} association membership \textbullet{} contact found \\
    Q4 Switching & Uses software Heidi integrates with \\
    Q5 Act now & Remote-reading meters mentioned \textbullet{} current metering provider \\
    Other ideas & District \\
    \bottomrule
  \end{tabular}
  \vspace{0.6em}

  12 candidates. Next, the tests decide which ones actually score.
  \note{I deliberately listed more candidates than I needed. The point of the tests is to
  throw most of them out for a stated reason.}
\end{frame}

%====================================================================
\section{Testing the KPIs}

\begin{frame}{Step 3: the tests}
  \partlabel{Part 4 \textbullet{} Testing every candidate}
  \footnotesize
  \begin{tabular}{@{}p{3.4cm} c c c p{4.8cm}@{}}
    \toprule
    \textbf{Candidate} & \textbf{Correl.} & \textbf{Discrim.} & \textbf{Feasible} & \textbf{Result} \\
    \midrule
    Units managed & \yes & \yes & partly & \textbf{Critical} (weight 3) \\
    Buyer type & \yes & \yes & \yes & \textbf{Important} (2) \\
    Residential share & \yes & weak & \yes & \textbf{Important} (2) \\
    Heidi-compatible software & partly & \yes & rarely & \textbf{Useful} (1) \\
    Association, growth & unproven & \yes & only if present & \textbf{Bonus} (+3 each) \\
    Buildings, employees & same as units & & & Only to \emph{estimate} units \\
    Remote reading mentioned & \yes & \no & & Info column \\
    Current metering provider & unclear & & rarely & Info column \\
    Contact found & \no & & & Separate column \\
    District & \no & & & Dropped \\
    \bottomrule
  \end{tabular}
  \note{Two kinds of failure here. District fails correlation: nothing says a company in
  Spandau is a better customer than one in Mitte. Remote reading fails discrimination: almost
  no company mentions it, so it can't separate good leads from bad.}
\end{frame}

\begin{frame}{The hard calls, and why}
  \begin{itemize}
    \item \textbf{Residential share was demoted from Critical to Important.} It's relevant, but
          most Hausverwaltungen are mainly residential, so it barely separates them.
    \item \textbf{Buildings and employees aren't scored.} They measure the same thing as
          units. Scoring all three would count size three times.
    \item \textbf{Association and growth became bonus points.} We can only see them when
          they're present. ``Not found'' doesn't mean ``no''.
    \item \textbf{Contact found isn't part of the score.} Whether we found an email says
          nothing about how good a customer the company would be. It's actionability, so it
          gets its own column.
    \item \textbf{Municipal companies aren't excluded by hand.} They became a value of buyer
          type, scored on evidence.
  \end{itemize}
  \note{Each of these was a decision I could have made by feel. Writing down the reason forced
  me to make them consistent.}
\end{frame}

\begin{frame}{The final KPI set}
  \small
  \begin{tabular}{@{}p{3.3cm} p{2.0cm} p{2.6cm} p{5.0cm}@{}}
    \toprule
    \textbf{KPI} & \textbf{Weight} & \textbf{Segment} & \textbf{Scale (0--5)} \\
    \midrule
    Units managed & 3 & Firmographic & 5 bands from real published values \\
    Buyer type & 2 & Firmographic & Rental 5 \textbullet{} WEG / mixed 4 \textbullet{}
      coop / own 3 \textbullet{} municipal 2 \\
    Residential share & 2 & Firmographic & Residential 5 \textbullet{} mixed 3
      \textbullet{} commercial 1 \\
    Heidi software & 1 & Technographic & Integrated 5 \textbullet{} other 2 \\
    Association & +3 & Behavioral & Only adds \\
    Growth signal & +3 & Behavioral & Only adds \\
    \bottomrule
  \end{tabular}
  \vspace{0.6em}

  All three segments are covered: \textbf{what the company is}, \textbf{what it uses},
  \textbf{what it's doing}.
  \note{The segmentation check confirms nothing big is missing: we have size and type, the
  technology, and activity signals.}
\end{frame}

%====================================================================
\section{The evidence behind the KPIs}

\begin{frame}{Evidence: size and who signs}
  \partlabel{Part 5 \textbullet{} What backs each KPI}
  \begin{itemize}
    \item \textbf{Units managed (weight 3):} Heidi prices \emph{per housing unit per year}, so
          revenue scales directly with units \hfill \ev{} (Heidi blog)
    \item \textbf{Buyer type, rental manager = 5:} for rental property the landlord decides,
          or the manager if authorised, so one relationship covers many owners
          \hfill \ev{} (Heidi blog)
    \item \textbf{Buyer type, WEG manager = 4:} the owners choose the metering provider by
          resolution, usually on the manager's proposal \hfill \ev{} (\S\,19 WEG)
    \item \textbf{Buyer type, municipal = 2:} WBM and HOWOGE run metering through their own
          subsidiaries \hfill \ev{} (job ad, tender)
  \end{itemize}
  \vspace{0.4em}
  \small A detail I caught: Heidi's own blog cites \S\,21 WEG. Since the 2020 reform the
  correct basis is \textbf{\S\,19 WEG}.
  \note{The WEG point is why WEG managers score 4 and not 5: there's an extra approval step,
  the owners' vote.}
\end{frame}

\begin{frame}{Evidence: switching and timing}
  \begin{itemize}
    \item \textbf{Software:} Heidi integrates with Domus, Immoware24, Wodis Sigma and
          Immocloud \hfill \ev{} (Heidi blog)
    \item \textbf{Why now:}
      \begin{itemize}
        \item Existing meters that can't be read remotely must be retrofitted or replaced
              \textbf{by 31 Dec 2026} \hfill \ev{} (HeizkostenV \S\,5 Abs.\,3)
        \item Otherwise tenants may cut their heating-cost share by \textbf{3\%}
              \hfill \ev{} (\S\,12)
        \item Metering contracts often run \textbf{up to 10 years} \hfill \ev{} (Heidi blog)
      \end{itemize}
    \item Q5 isn't scored: the deadline applies to \emph{every} company, so it can't rank them.
          It's the \textbf{pitch}, not a KPI.
  \end{itemize}
  \note{The legal text was checked on gesetze-im-internet.de, not just in vendor blogs, because
  metering companies have an interest in the deadline.}
\end{frame}

\begin{frame}{What is still an assumption}
  Stated openly instead of presented as fact:
  \begin{itemize}
    \item \as{} Residential portfolios need more heating-cost billing. It's a proxy: the
          HeizkostenV also covers commercial buildings.
    \item \as{} Cooperatives and private landlords score 3, because they're one decision-maker
          but may have central contracts.
    \item \as{} Software integration raises the win rate.
    \item \as{} Association membership and growth signal a better customer.
    \item \as{} Managers reconsider the provider at retrofit time.
  \end{itemize}
  \vspace{0.6em}
  \textbf{Next step:} test these against Heidi's real win and loss data, and replace the weights.
  \note{A scoring model with hidden assumptions can't be improved. One with labelled
  assumptions tells Heidi exactly what to measure next.}
\end{frame}

%====================================================================
\section{Turning KPIs into a score}

\begin{frame}{Problem 1: missing data}
  \partlabel{Part 6 \textbullet{} From KPIs to one number}
  Most websites don't publish everything. What should ``unknown'' count as?
  \vspace{0.4em}

  \small
  \begin{tabular}{@{}p{5.2cm} p{3.6cm} p{3.6cm}@{}}
    \toprule
    & \textbf{Unknown = 0} & \textbf{Unknown left out} \\
    \midrule
    A: small, publishes everything & 66 & 66 \\
    B: large, publishes only its size & \textcolor{bad}{35} & 94 + \textbf{Low confidence} \\
    \bottomrule
  \end{tabular}
  \vspace{0.6em}

  \normalsize
  \begin{itemize}
    \item Scoring unknown as 0 punishes companies for having a short website
    \item So unknown KPIs are \textbf{left out of the average}, and a \textbf{Data
          completeness} column and a \textbf{confidence flag} show how much is known
  \end{itemize}
  \note{Company B is probably the better lead, because size is the critical KPI. Zero-scoring
  would hide it at 35. Leaving unknowns out puts it on top, but flagged, so sales knows to
  check it.}
\end{frame}

\begin{frame}{Problem 2: a bonus that lowered the score}
  First version: association membership scored 4/5 \emph{inside} the average.
  \vspace{0.4em}
  \begin{itemize}
    \item A company scoring 5 on everything: \textbf{100}
    \item The same company, plus a found membership (4): \textcolor{bad}{\textbf{98}}
  \end{itemize}
  Finding a \emph{positive} signal pulled the score \emph{down}.
  \vspace{0.6em}

  \textbf{Fix:} bonus signals sit outside the average and can only add.
  \[
  \text{Lead score} = \underbrace{\text{Base} \times 18.8}_{\text{max } 94} +
  \underbrace{\text{bonus}}_{\text{max } 6}
  \]
  \note{18.8 is 94 divided by 5. The base is scaled to 94, not 100, so the bonus still breaks
  ties at the top.}
\end{frame}

\begin{frame}{Problem 3: where do the size bands come from?}
  \begin{itemize}
    \item No industry standard says what a ``large'' Berlin Hausverwaltung is
    \item So the bands are set from the \textbf{178 companies that publish their units}, as
          five equal groups:
  \end{itemize}
  \begin{center}
  \begin{tabular}{ccccc}
    \toprule
    1 & 2 & 3 & 4 & 5 \\
    \midrule
    $<$800 & 800--1,400 & 1,400--3,000 & 3,000--6,640 & $>$6,640 \\
    \bottomrule
  \end{tabular}
  \end{center}
  \begin{itemize}
    \item When units aren't published, they're \textbf{estimated} from our own data: 150 units
          per employee (64 companies) or 29 per building (38 companies)
    \item Estimates are marked \emph{estimated}, and they don't shape the bands
  \end{itemize}
  \note{Calibrating from our own data means the scale fits the Berlin market, not an
  arbitrary rule of thumb.}
\end{frame}

%====================================================================
\section{Review}

\begin{frame}{How review changed the KPIs}
  \partlabel{Part 7 \textbullet{} Iteration}
  The KPI set was reviewed twice before anything was built. What changed:
  \begin{itemize}
    \item Residential share: \textbf{Critical \textrightarrow{} Important} (weak discrimination)
    \item Remote reading: \textbf{KPI \textrightarrow{} info column} (almost nobody mentions it;
          in the end, 20 of 1,563 sites)
    \item Association: \textbf{inside the average \textrightarrow{} bonus} (the 98 bug)
    \item WEG basis: \textbf{\S\,21 \textrightarrow{} \S\,19 WEG}
    \item Municipal companies: \textbf{exclusion by hand \textrightarrow{} buyer type value}
    \item Unit estimation: \textbf{rule of thumb \textrightarrow{} calibrated on our own data}
    \item Lead score and contact status: \textbf{split} into two columns
  \end{itemize}
  \note{Having a written method made the review productive. Each critique pointed at a
  specific rule, and each fix was a rule change, not a rewrite.}
\end{frame}

%====================================================================
\section{The agent}

\begin{frame}{From KPIs to the agent}
  \partlabel{Part 8 \textbullet{} Filling the KPIs automatically}
  Each KPI definition became an instruction for the agent. For example:
  \vspace{0.4em}

  \small
  \begin{block}{Buyer type, as the agent receives it}
    \emph{third\_party\_rental:} manages property for other owners, with rental management as
    the stated or clearly dominant focus. \emph{third\_party\_weg:} \dots{} with
    WEG-Verwaltung as the focus. \emph{not\_property\_manager:} only brokerage, caretaking,
    construction or law. A broker that also manages property is NOT this.
  \end{block}
  \normalsize
  \begin{itemize}
    \item Same definitions for the agent, the scoring and the hand check
    \item Each value must come with a \textbf{source URL and a word-for-word quote}, or it's
          ``unknown''
  \end{itemize}
  \note{Because the KPIs were precisely defined, they could be handed to an AI. A vague KPI like
  ``good fit'' can't be automated or checked.}
\end{frame}

\begin{frame}{The pipeline: AI for judgement, code for exactness}
  \begin{enumerate}
    \item \textbf{Collect} (code): Google Places, with Berlin split into map cells, plus
          OpenStreetMap
    \item \textbf{Merge} (code): duplicates by website domain, then by name
    \item \textbf{Fetch} (code): up to 6 pages per site, robots.txt respected
    \item \textbf{Research} (AI, Claude Haiku 5.5): reads the pages, may open 4 more, submits
          facts with quotes
    \item \textbf{Verify} (code): every quote checked against the page; a failure becomes
          unknown
    \item \textbf{Score \& export} (code): the formula, written to Excel
  \end{enumerate}
  \vspace{0.4em}
  One command, no manual steps. Budget stop, automatic quality gate after 25 companies,
  checkpoint after every company.
  \note{The AI never calculates a score. It reads and extracts; code checks and scores.}
\end{frame}

%====================================================================
\section{Does it work?}

\begin{frame}{Validation}
  \partlabel{Part 9 \textbullet{} Can the numbers be trusted?}
  \begin{tabular}{@{}p{8.0cm} p{4.8cm}@{}}
    \toprule
    \textbf{Check} & \textbf{Result} \\
    \midrule
    Claims confirmed on the cited page (code) & 3,945 / 4,127 (96\%) \\
    Accuracy vs. hand-checked websites & 41 / 45 fields (91\%) \\
    Haiku vs. Sonnet on the same 5 companies & 21 vs. 22 of 25 \\
    Cost per company: Haiku vs. Sonnet & \$0.002 vs. \$0.033 \\
    \bottomrule
  \end{tabular}
  \vspace{0.6em}
  \begin{itemize}
    \item Every miss was a cautious ``unknown'' or a borderline call, never an invented value
    \item The agent even found a figure the hand check had missed
    \item Sonnet is barely better at 15 times the cost, so Haiku runs the list
  \end{itemize}
  \note{Two different checks: the 96\% proves the quotes are real; the 91\% proves the answers
  are right.}
\end{frame}

\begin{frame}{Results: 2,080 Berlin companies for \$2.93}
  \begin{columns}[T]
    \column{0.5\textwidth}
    \begin{itemize}
      \item \textbf{242} scored, confidence OK
      \item \textbf{904} scored, confidence Low
      \item \textbf{417} excluded: not property managers
      \item \textbf{219} sites dead or blocked
      \item \textbf{517} without own website
    \end{itemize}
    \column{0.5\textwidth}
    \begin{itemize}
      \item \textbf{1,098} named contacts
      \item \textbf{88} with a personal email
      \item \textbf{1,070} company emails
      \item Every value has its source and quote in the sheet
    \end{itemize}
  \end{columns}
  \note{The Excel file is sorted by lead score, with confidence and contact status next to it.}
\end{frame}

%====================================================================
\section{What the data said about the KPIs}

\begin{frame}{What the data said back about the KPIs}
  \partlabel{Part 10 \textbullet{} Lessons}
  \begin{itemize}
    \item \textbf{The most important KPI is the hardest to find:} only 178 companies publish
          their units. That's why the confidence flag and unit estimation matter.
    \item \textbf{Buyer type was worth defining carefully:} it caught 417 brokers and other
          non-managers that Google listed as property management.
    \item \textbf{Software got the right low weight:} only 82 companies name theirs.
    \item \textbf{Municipal = 2 was confirmed:} the agent found that degewo and STADT UND LAND
          also run metering in-house.
    \item \textbf{Dropping remote reading was right:} only 20 sites mention it.
  \end{itemize}
  \note{The data validated most of the KPI decisions made before any data existed. That's the
  benefit of deciding the method first.}
\end{frame}

\begin{frame}{Next steps}
  \begin{itemize}
    \item \textbf{Turn assumptions into evidence:} fit the weights on Heidi's real wins and
          losses
    \item \textbf{Close the size gap:} add sources for units (press, directories, registers)
    \item \textbf{Re-check the top leads} with a stronger model
    \item \textbf{Cover the 517 companies without a website}
    \item \textbf{Run it weekly:} only new companies are researched, at about \$0.002 each
  \end{itemize}
  \vspace{1em}
  Code, method and write-up: \textbf{github.com/6hammad9/heidi-lead-pipeline}
  \note{The method is built to be improved: every weight is in one config file, and every
  assumption is labelled, so Heidi's data can replace them one by one.}
\end{frame}

\end{document}
